\documentclass[11pt]{article}
\usepackage[T1]{fontenc}
\usepackage{lmodern,amsmath,amssymb,amsthm,mathtools,booktabs,microtype}
\usepackage[margin=1.05in]{geometry}
\usepackage{xcolor}
\usepackage[colorlinks=true,linkcolor=blue!55!black,urlcolor=blue!55!black,citecolor=blue!55!black]{hyperref}
\setlength{\emergencystretch}{3em}
\newtheorem{theorem}{Theorem}[section]
\newtheorem{proposition}[theorem]{Proposition}
\newtheorem{corollary}[theorem]{Corollary}
\newtheorem{lemma}[theorem]{Lemma}
\newtheorem{example}[theorem]{Example}
\theoremstyle{remark}\newtheorem{remark}[theorem]{Remark}
\newcommand{\R}{\mathbb R}\newcommand{\C}{\mathbb C}\newcommand{\supp}{\operatorname{supp}}
\newcommand{\ord}{\operatorname{ord}}\newcommand{\diam}{\operatorname{diam}}
\newcommand{\Rea}{\operatorname{Re}}\newcommand{\Ima}{\operatorname{Im}}\newcommand{\dd}{\,\mathrm d}
\newcommand{\e}{\mathrm e}\newcommand{\norm}[1]{\left\|#1\right\|}
\title{Sharp Positive Zero-Cancelling Filters\\for Paired and Perturbed Spectra}
\author{Research note for the ProveIt GRH project}
\date{27 November 2026}
\begin{document}
\maketitle
\begin{abstract}
A positive average of dilations can cancel specified zero terms in a prime explicit formula. We determine its exact minimum logarithmic width for arbitrary finite multisets of complex exponents paired by reflection about a horizontal line, including repeated vanishing constraints. The optimum is characterized by coefficient positivity of a single real polynomial. When all coefficients are strictly positive, arbitrary small perturbations of the exponents preserve a unique globally optimal atomic measure, whose nodes, weights and width are real analytic in the exponents. We give a sensitivity formula, a sharp Wasserstein stability exponent and the first-order cost of a shorter interval. As applications, we solve the equal-jet two-ordinate problem, obtain exact three-frequency formulas beyond the consecutive-harmonic case, and certify the optimum for a rational box containing the first three zeta ordinates: its width lies in $(0.46994519,0.46994524)$. The extremal theorems are unconditional; the use of the supplied Dirichlet GRH theorem in prime explicit formulas is made explicit. The interval certificate uses only integer arithmetic.
\end{abstract}
\tableofcontents

\section{Source, scope and contribution}
The GitHub snapshot inspected is \texttt{VladimirReshetnikov/ProveIt}, main ref
\begin{center}\small\texttt{a4404ca0247c31b8760589015b7eb6ab4efb6f92}.\end{center}
We inspected the vendoring README, the vendored README, \texttt{src/Statements.lean}, \texttt{src/Main.lean}, the merged article README, its introduction and source section, and the relevant reconstruction section, particularly \texttt{rec:sec:dilation}. We also surveyed headings in the Li, zero and direct-consequences sections. This is a targeted reading rather than a line-by-line audit of the 723-page article or of the approximately 300,000 lines of Lean code.

The vendoring README pins the readable edition of \texttt{lsernau/grh-lean} at the commit
\begin{center}\small\texttt{72fead102f663d5c7a407964a13dc0cf9d977207}.\end{center}
Its five main declarations are
\begin{quote}\small\ttfamily
riemannHypothesis\\
riemannZeta\_eq\_zero\_re\\
LFunction\_ne\_zero\_of\_half\_lt\_re\\
LFunction\_eq\_zero\_re\\
LFunction\_eq\_zero\_re\_of\_isPrimitive
\end{quote} In particular, for every Dirichlet character $\chi$, a zero of $L(s,\chi)$ in $0<\Rea s<1$ has real part $1/2$. The source README reports standard Lean axioms only and independent checker acceptance. The vendored package is not checked by ProveIt's root build; we did not replay it here. We use the exported assertion as input, denoted DGRH, without claiming a fresh verification. It does not imply RH for arbitrary Artin or Hecke $L$-functions, zero simplicity, or linear independence of ordinates. The article itself expressly disclaims independent checking of its merged proofs \cite{vendor,article}.

The arithmetic mechanism in the article's source section compares covolume bounds, then uses a one-sided smoothed prime remainder to detect zeros. The reconstruction section develops a complementary positivity-preserving operation: average dilations to annihilate selected spectral terms. Its theorem \texttt{rec:thm:hblock} establishes the exact width for consecutive harmonics by a keyhole-residue certificate; \texttt{rec:thm:twofreq} gives the two-frequency optimum for simple vanishing, and \texttt{rec:thm:cusp} describes resonant transitions. The open questions at the end of that section explicitly ask about three noncommensurate frequencies and derivative constraints at two frequencies. Here we generalize the keyhole construction to arbitrary reflection-paired exponents, including a coefficient-sign criterion, and then prove persistence of the optimum under independent perturbations. The second step is important: the exact optimizer is not restricted to a grid, and a perturbation need not retain any reflection relation. We also give a verified three-frequency application beyond the small-perturbation existence statement.

Minimum-duration positive impulse shaping and derivative robustness have a substantial earlier literature, including \cite{lau,pao}. Thus the problem formulation, convolution, and optimality certificates are not themselves claims of originality. The particular general reflection/perturbation theorem and the uniform three-ordinate interval certificate are extensions relative to the inspected article; a limited literature search did not establish priority. No claim of resolution of a famous open number-theoretic conjecture is made.

\section{The extremal problem and a global exact theorem}
For a compactly supported Borel probability measure $\mu$ on $\R$ put
\[
 F_\mu(s)=\int\e^{su}\dd\mu(u),\qquad
 D(\boldsymbol{s},\boldsymbol{m})=\inf_\mu\diam\supp\mu,
\]
where the infimum is over measures such that $F_\mu^{(k)}(s_j)=0$ for $0\le k<m_j$, for distinct $s_j$ in the open upper half-plane and integers $m_j\ge1$. Let $n=\sum_j m_j$. Translation multiplies $F_\mu$ by an everywhere nonzero exponential, hence preserves the constraints and the width. The feasible set is nonempty: at a single $s=\beta+i\omega$, the two masses at $0,\pi/\omega$ in ratio $\e^{\beta\pi/\omega}:1$ give a simple zero; convolve the requisite number of these two-point measures for each constraint. The infimum is attained, by translating minimizing measures to a fixed compact interval, taking a weak limit and using continuity of the test functions and the support bound. These facts require no RH.

\begin{theorem}[Reflection-paired spectra]\label{thm:base}
Fix $h>0$. Suppose $0<\Ima s_j<2\pi/h$ and the multiset of targets, with multiplicities $m_j$, is invariant under
\[
 s\longmapsto \bar s+2\pi i/h.
\]
Define the real monic polynomial
\[
 P(z)=\prod_j(z-\e^{hs_j})^{m_j}=\sum_{k=0}^n p_k z^k.
\]
Then $P(0)>0$, $P(1)>0$, and
\[
 \boxed{D(\boldsymbol{s},\boldsymbol{m})\ge nh.}
\]
Equality holds if and only if every $p_k\ge0$. In the equality case the unique minimizer up to translation is
\[
 \boxed{\displaystyle \mu_*=\frac{1}{P(1)}\sum_{k=0}^n p_k\delta_{kh}.}
\]
If at least one coefficient is negative, the optimum is strictly larger than $nh$. No assumption of commensurability or a common real part is required.
\end{theorem}
\begin{proof}
Set $q_j=-\e^{hs_j}$, $Q(z)=(-1)^nP(-z)=\prod_j(z-q_j)^{m_j}$. The $q_j$ lie in the slit plane $\C\setminus(-\infty,0]$, and their multiset is closed under conjugation. The real members are positive; other members form conjugate pairs. Consequently $Q(-t)=(-1)^n|Q(-t)|$ for $t\ge0$, $P(0)>0$ and $P(1)>0$. Distinct targets give distinct roots, since their imaginary parts lie in an open interval of length $2\pi/h$.

For real $x$ use the principal logarithm and define
\[
 E(x)=\sum_j\operatorname{Res}_{z=q_j}\frac{z^{x-1}}{Q(z)},\qquad
 G(u)=(-1)^{n+1}\sin(\pi u/h)E(u/h).
\]
The residue formula, including at repeated poles, makes $E(x)$ a finite linear combination of $x^k\e^{x\log q_j}$ with $0\le k<m_j$. It is real for real $x$. Because $\log q_j=hs_j-i\pi$ on the chosen branch, multiplying by $\sin(\pi u/h)$ produces terms at exponents $s_j$ and $s_j-2\pi i/h=\bar s_{j'}$ for the reflected partner $j'$. Therefore $G$ is a real linear combination of the real and imaginary parts of $u^k\e^{s_j u}$, $k<m_j$. In particular $\int G\dd\mu=0$ for every feasible measure.

For $0<x<n$, apply the residue theorem to a keyhole contour about the negative real axis. The large and small circular contributions tend to zero, since the powers are $R^{x-n}$ and $\epsilon^x$. The jump across the cut is $2i\sin(\pi(x-1))t^{x-1}/Q(-t)$, with upper bank traversed toward zero. Hence
\[
 E(x)=\frac{(-1)^{n+1}\sin(\pi x)}{\pi}
       \int_0^\infty\frac{t^{x-1}}{|Q(-t)|}\dd t,
\]
and thus
\begin{equation}\label{eq:certificate}
 G(hx)=\frac{\sin^2(\pi x)}{\pi}
       \int_0^\infty\frac{t^{x-1}}{|Q(-t)|}\dd t\qquad(0<x<n).
\end{equation}
The integral is strictly positive and finite. The endpoint asymptotics are obtained by splitting at $1$: it is $1/(|Q(0)|x)+O(1)$ as $x\downarrow0$, and $1/(n-x)+O(1)$ as $x\uparrow n$. In particular $G$ has simple endpoint zeros, double interior zeros, and
\begin{equation}\label{eq:slopes}
 G'(0)=\frac{\pi}{h|Q(0)|},\qquad G'(nh)=-\frac{\pi}{h}.
\end{equation}
It is positive elsewhere on $[0,nh]$.

Any feasible measure of width at most $nh$ may be translated into this interval. The zero integral of the nonnegative $G$ forces it onto the nodes $0,h,\ldots,nh$. Its transform is $A(\e^{hs})$ for a real polynomial $A$ of degree at most $n$ with nonnegative coefficients and $A(1)=1$. The map $s\mapsto\e^{hs}$ has nonzero derivative, so all the prescribed zeros of $A$ have the required multiplicities. Thus $A=P/P(1)$, including its nonzero constant and leading coefficients. It follows that any such feasible measure has width exactly $nh$, and exists precisely when every coefficient is nonnegative. Attainment of the global infimum, shown above, gives strict inequality in the remaining case.
\end{proof}

\begin{corollary}[Two ordinates with equal jet orders]\label{cor:pair}
Let $0<\alpha<\gamma\le3\alpha$, $\beta\in\R$, and $m\ge1$. Requiring a zero of order at least $m$ at both $\beta+i\alpha$ and $\beta+i\gamma$ has exact optimal width
\[
 D=\frac{4\pi m}{\alpha+\gamma}.
\]
With $h=2\pi/(\alpha+\gamma)$, $a=\e^{h\beta}$ and $\theta=\pi-h\alpha$, the optimal weights at $0,h,\ldots,2mh$ are the coefficients of
\[
 \frac{(z^2+2a\cos\theta\,z+a^2)^m}{(1+2a\cos\theta+a^2)^m}.
\]
\end{corollary}
\begin{proof}
The exponents are reflected partners and $0<\theta\le\pi/2$. All coefficients are nonnegative. Apply Theorem~\ref{thm:base}. The same formula is the $m$-fold convolution of the simple optimum, but the theorem proves that no other positive measure, including a continuous one, can shorten it. At the boundary ratio $3$, alternate coefficients vanish and the support is smaller than $2m+1$.
\end{proof}

\begin{corollary}[A nonconsecutive three-frequency band]\label{cor:triple}
Let $0<\alpha<\gamma\le5\alpha$, $b=(\alpha+\gamma)/2$, and $\beta\in\R$. For simple vanishing at $\beta+i\alpha,\beta+ib,\beta+i\gamma$, the exact width is $6\pi/(\alpha+\gamma)$. The polynomial, with $h,a,\theta$ as in Corollary~\ref{cor:pair}, is
\[
 P(z)=(z+a)(z^2+2a\cos\theta\,z+a^2)
 =z^3+a(1+2\cos\theta)z^2+a^2(1+2\cos\theta)z+a^3.
\]
For $\gamma<5\alpha$ every coefficient is positive; at $\gamma=5\alpha$ only the endpoints remain. For $\gamma>5\alpha$ this width is unattainable. Raising this polynomial to the $m$th power gives the exact width $6\pi m/(\alpha+\gamma)$ for equal order $m$ at all three, in the same band.
\end{corollary}
\begin{proof}
Here the central target is a fixed point of reflection, and $1+2\cos\theta\ge0$ is equivalent to $\gamma/\alpha\le5$. Theorem~\ref{thm:base} applies. If the ratio exceeds $5$, the coefficient of degree $3m-1$ in $P^m$ is negative.
\end{proof}
For example, choose $\cos\theta=-1/3$, $\beta=0$. The exact weights for the simple triple are $(3,1,1,3)/8$. The ratio $\gamma/\alpha$ exceeds $3$, so the optimal two-frequency formula of Corollary~\ref{cor:pair} alone does not establish this result. Arbitrary real angles in the allowed ranges are permitted, unlike the consecutive-harmonic case.

\section{Independent perturbations: global uniqueness and analytic dependence}
Let $\sigma=(s_1,\ldots,s_r)$, with multiplicities fixed, and write $Y_\sigma(u)\in\R^{2n}$ for the real and imaginary parts of $u^k\e^{s_j u}$. The parameters are the $2r$ real coordinates of the targets. At a strictly positive base point of Theorem~\ref{thm:base}, denote the normalized weights by $w_i$, and nodes by $u_i=ih$.

\begin{theorem}[Analytic continuation of the optimum]\label{thm:perturb}
If all coefficients of $P$ at the base point are strictly positive, there is an open neighborhood of the base targets in which the unique global minimum-width measure, with left endpoint fixed at $0$, has the form
\[
 \mu_\sigma=\sum_{i=0}^n w_i(\sigma)\delta_{u_i(\sigma)},\qquad
 0=u_0<u_1<\cdots<u_n,\quad w_i>0.
\]
All weights and nodes, and therefore the optimal width $L(\sigma)=u_n(\sigma)$, are real analytic. The targets may move independently in both real and imaginary directions. No reflection symmetry is needed at the perturbed point.
\end{theorem}
\begin{proof}
The finite system for the $2n+1$ real unknowns $(w_0,\ldots,w_n,u_1,\ldots,u_n)$ is
\begin{equation}\label{eq:system}
 \sum_i w_i=1,\qquad \sum_i w_iY_\sigma(u_i)=0.
\end{equation}
Its Jacobian in the unknowns at the base point is invertible. To see this, suppose an infinitesimal variation belongs to its kernel. The variation of the transform is
\[
 \delta F(s)=A(\e^{hs})+s B(\e^{hs}),\quad
 A(z)=\sum_i\delta w_i z^i,\quad B(z)=\sum_i w_i\delta u_i z^i.
\]
The polynomials are real, $A(1)=0$, and $B(0)=0$ because $u_0$ is fixed. All target jets of $\delta F$ vanish. By conjugating the partner target's equations, $\delta F(s-2\pi i/h)$ also vanishes to the required order at every target $s_j$. Subtracting these two functions cancels $A$ and yields $(2\pi i/h)B(\e^{hs})$. Thus $B$ is divisible by $P$. Since $\deg B\le n$ and $B(0)=0\ne P(0)$, $B=0$. The same argument for the remaining jets gives $A$ divisible by $P$; $A(1)=0\ne P(1)$ implies $A=0$. Positivity of all $w_i$ forces every node variation to vanish. The real analytic implicit function theorem gives an analytic solution of \eqref{eq:system} near the base, still with positive weights and strictly ordered nodes, and with invertible Jacobian.

We must prove global, not merely local, optimality. In $\R^{2n}$ consider the columns $Y(u_i)$ for $0\le i\le n$ and $Y'(u_i)$ for $1\le i<n$. Their span has dimension $2n-1$: the weight relation in \eqref{eq:system} gives at least one dependence among the $2n$ columns; adding the last derivative column gives rank $2n$ by the full Jacobian. Choose a nonzero normal $v_\sigma$ to this hyperplane and normalize by $v_\sigma\cdot Y'(u_n)=-1$. The normalization is possible by full rank, and the normal depends real analytically on the parameters (solve a nonsingular minor system). The function
\[
 g_\sigma(u)=v_\sigma\cdot Y_\sigma(u)
\]
vanishes at every node, with vanishing first derivative at the interior nodes. At the base point it is a positive multiple of $G$ in \eqref{eq:certificate}, since the base normal space is one dimensional and $G'(nh)<0$. Consequently its left endpoint derivative is positive, its interior second derivatives are positive, and its right endpoint derivative is negative. These properties persist under perturbation, by continuity. Positivity off small neighborhoods of the nodes persists by compactness (for example rescale the interval by its right endpoint); Taylor's theorem and the strict derivative signs give positivity in the node neighborhoods. Shrinking the parameter neighborhood if needed, $g_\sigma\ge0$ on $[0,L(\sigma)]$, with exactly the listed zero set.

Any feasible measure of width at most $L(\sigma)$ can be translated into this interval, and is forced onto the nodes by the integral of $g_\sigma$. The weight columns in the full Jacobian are linearly independent after adding the mass row, so \eqref{eq:system} fixes the weights uniquely on those nodes. Both endpoint weights are positive. Hence the analytic candidate is the unique global minimizer up to translation.
\end{proof}

\begin{proposition}[Sensitivity of the minimum]\label{prop:sense}
At every point in the neighborhood of Theorem~\ref{thm:perturb}, for an infinitesimal parameter change $\dot\sigma$,
\begin{equation}\label{eq:sense}
 \dot L=-\frac{\sum_{i=0}^n w_i v_\sigma\cdot (\partial_\sigma Y_\sigma)(u_i)[\dot\sigma]}
                    {w_n g_\sigma'(L)}.
\end{equation}
\end{proposition}
\begin{proof}
Differentiate the moment equation in \eqref{eq:system}, then dot with the normal $v_\sigma$. All weight terms and interior node terms vanish; $u_0$ is fixed, leaving the displayed formula. The denominator is nonzero.
\end{proof}

\begin{example}[Three independently detuned harmonics]\label{ex:detune}
For Fourier frequencies $(1+\epsilon_1,2+\epsilon_2,3+\epsilon_3)$ near $(1,2,3)$, the minimum is
\[
 L=\frac{3\pi}{2}-\frac\pi4(\epsilon_1+\epsilon_2+\epsilon_3)+O(\norm{\epsilon}^2).
\]
Indeed at the base $h=\pi/2$ and $P=1+z+z^2+z^3$. By uniqueness and reflection, for pure imaginary targets the centered minimizer is even, with nodes $\pm t,\pm T$ and total inner mass $w$. At the base $(t,T,w)=(\pi/4,3\pi/4,1/2)$. Its three cosine equations are
\begin{equation}\label{eq:cos}
 w\cos(\lambda_j t)+(1-w)\cos(\lambda_j T)=0\quad(j=1,2,3).
\end{equation}
Differentiating the first and third and adding eliminates $\dot w$, giving $\dot t+\dot T=-\pi(\epsilon_1+\epsilon_3)/4$; the second gives $\dot T-\dot t=-\pi\epsilon_2/4$. Since $L=2T$, the stated expansion follows. In particular nearby rationally independent triples have a unique analytic global optimizer. This asserts a neighborhood, not an explicit universal size for it.
\end{example}
\begin{example}[An exact slice off the reflection locus]
For $b$ sufficiently close to $2$, the frequencies $(1,b,3)$ have exact optimal width $\pi+\pi/b$, even for irrational $b$. The convolution of the two centered equal-mass filters at $1$ and $b$ annihilates all three frequencies: the filter at $1$ also annihilates $3$. Its four nodes are $\pm(\pi/2-\pi/(2b))$, $\pm(\pi/2+\pi/(2b))$, each of mass $1/4$. This candidate equals the base optimizer at $b=2$ and solves \eqref{eq:system} analytically, so the implicit solution is this measure and Theorem~\ref{thm:perturb} gives global optimality. For $b\ne2$ it does not come from the reflection pairing used at the base.
\end{example}

\section{Stability and the exact cost of shortening}
Fix a point of Theorem~\ref{thm:perturb}, including the strictly positive base, and write $Y,g=v\cdot Y,L,u_i,w_i$ for its data. Use any norm on $\R^{2n}$, its dual $\norm{\cdot}_*$, and $\varepsilon(\nu)=\norm{\int Y\dd\nu}$. All estimates in this section hold for fixed targets; constants need not be uniform near coefficient zeros or colliding targets.

\begin{theorem}[Sharp transport stability]\label{thm:stable}
For every probability measure $\nu$ on $[0,L]$ with $\varepsilon(\nu)\le1$,
\[
 W_1(\nu,\mu_*)\le C\varepsilon(\nu)^{1/2}\quad(n\ge2).
\]
If $n=1$, the exponent is $1$. Both exponents are optimal.
\end{theorem}
\begin{proof}
For $n\ge2$, positivity and the simple and double zeros imply $g(u)\ge c\,\operatorname{dist}(u,\{u_i\})^2$ throughout $[0,L]$ for some $c>0$. The integral of $g$ is at most $\norm{v}_*\varepsilon$. Moving each point of $\nu$ to its nearest node costs $O(\sqrt\varepsilon)$ in $W_1$ by Cauchy--Schwarz. The move changes the moments by at most a Lipschitz constant times that cost. The injective mass-plus-moment matrix of the nodes has a bounded inverse on its image, so the rounded weights differ from the optimal ones in total variation by $O(\sqrt\varepsilon)$. Moving between nodes costs at most $L$ times that variation. For $n=1$, both zeros are simple and the lower bound on $g$ is linear in the distance, giving the linear estimate. For sharpness when $n\ge2$, replace one interior atom by two half-atoms displaced by $\pm a$. The moment change is $O(a^2)$ by Taylor's theorem, but the transport distance is at least its mass times $a$ (use the distance-to-nodes test function, which is $1$-Lipschitz). When $n=1$, move one endpoint mass inward by $a$ to obtain a transport cost proportional to $a$ and moment error $O(a)$.
\end{proof}

\begin{theorem}[Anchored shortening law]\label{thm:short}
For $d>0$ small, let $M(d)$ be the minimum of $\varepsilon(\nu)$ over probability measures on $[0,L-d]$. Then
\[
 \boxed{M(d)=\frac{w_n[-g'(L)]}{\norm{v}_*}\,d+O(d^2).}
\]
\end{theorem}
\begin{proof}
Set $a=-g'(L)>0$ and $K=w_na$. Since the $Y(u_i)$ have rank $n$ with unique relation proportional to $(w_i)$, there is a real linear combination $h_0=z\cdot Y$ with values $K$ at every node except the last, and $K-a$ at the last. Indeed these values have weighted sum zero. Near the interior nodes the quadratic lower bound for $g$ and the Lipschitz bound for $h_0$ imply $g+dh_0\ge dK-O(d^2)$. Near zero the derivative is strictly positive for small $d$, so the same bound follows from the value at zero. Near the last endpoint the derivative is strictly negative, so the minimum on the shortened interval is attained at $L-d$, with value $dK+O(d^2)$. Off these neighborhoods $g$ has a positive lower bound. Integrating gives
\[
 (\norm{v}_*+d\norm{z}_*)\varepsilon(\nu)\ge dK-O(d^2).
\]
For the upper bound, choose $r\in\R^{2n}$ with $v\cdot r=K$ and $\norm r=K/\norm v_*$ (existence follows from compactness of the unit ball and duality). Move the last endpoint by $-d$. Its first-order moment change is $-dw_nY'(L)$, having scalar product $dK$ with $v$. Vary the weights subject to total mass one and the $n-1$ interior nodes. Their first-order moment changes span $v^\perp$: they lie in that hyperplane, and injectivity of the full Jacobian on these $2n-1$ free coordinates gives equality of dimensions. Hence choose variations proportional to $d$ so the total moment is $dr+O(d^2)$. Weights and ordering remain valid for small $d$, and the support is inside $[0,L-d]$. This establishes the reverse inequality.
\end{proof}
At the paired base, \eqref{eq:slopes} gives $w_n[-G'(L)]=\pi/(hP(1))$. Thus the first-order constant is explicit once the coefficient vector of the residue certificate in the test basis and the norm are specified. An anchored interval is used here to avoid any ambiguity caused by translation of derivative constraints.

\section{A certified optimum for three specified ordinates}\label{sec:cert}
The preceding perturbation theorem alone does not certify how far a concrete triple can move. Here we give a separate finite proof, uniform over the following closed rational intervals:
\begin{equation}\label{eq:lambda}
\begin{split}
 \lambda_1&\in[14.1347251417,14.1347251418],\\
 \lambda_2&\in[21.0220396387,21.0220396388],\\
 \lambda_3&\in[25.0108575801,25.0108575802].
\end{split}
\end{equation}
These contain the first three positive zeta ordinates according to the documented FLINT consecutive-zero calculation \cite{flint}. Our certificate does not itself verify the location or indexing of the zeta zeros; it certifies the extremal result for \emph{every} triple in \eqref{eq:lambda}. The input zero-location data are thus clearly separated from the optimization proof.

\begin{theorem}[Three-ordinate certificate]\label{thm:actual}
For every triple in \eqref{eq:lambda}, the simple Fourier cancellation problem has a unique global minimizer up to translation. It is an even measure in centered coordinates with nodes $\pm t,\pm T$, total inner mass $w$, and equal masses within each pair. The parameters solve \eqref{eq:cos} uniquely in the closed box of radius $10^{-8}$ in the maximum norm about
\[
 (0.0794453460700645,\ 0.2349726063584985,\ 0.6943858328459642).
\]
Its minimum width satisfies
\[
 \boxed{0.46994519<D<0.46994524.}
\]
The finite checks needed in this theorem are implemented in \texttt{certify\_three.py}, using exact integer outward rounding; the executed transcript is \texttt{certificate.txt}.
\end{theorem}
\begin{proof}
Write $f(t,T,w)$ for the left sides of \eqref{eq:cos}, $x_0$ for the given center, $X$ for its box, and $r=10^{-8}$. The program constructs a rational matrix $A$ by inverting the matrix of interval midpoints of $Df(x_0)$; the latter intervals include all parameters in \eqref{eq:lambda}. Its rational interval checks establish, uniformly in the frequency box,
\[
 \sup_X\norm{I-A Df}_\infty<4.0\cdot10^{-6},\qquad
 \norm{A f(x_0)}_\infty<5.8\cdot10^{-12}.
\]
The contraction $x\mapsto x-Af(x)$ maps $X$ into itself by these bounds and is strictly contractive. $A$ is invertible (the bound on $I-A Df$ implies $A Df$ is invertible), so the unique fixed point satisfies $f=0$. The checks also give $0<t<T$, $0<w<1$.

For this solution put $a=(\cos(\lambda_jt))_j$, $b=(-\lambda_j\sin(\lambda_jt))_j$, and $c=\pm(a\times b)$ with sign chosen so $\sum c_j>0$. Set $g(u)=\sum_j c_j\cos(\lambda_ju)$. Algebraically $g(t)=g'(t)=0$ and $g(T)=0$ by \eqref{eq:cos}. For every solution in the box the program bounds $g''>0$ on $[0.075,0.084]$ and $g'<0$ on $[0.230,T_{\max}]$, with $t$ and $T$ in the corresponding neighborhoods. It proves $g>0$ on the remaining intervals $[0,0.075]$ and $[0.084,0.230]$ by subdividing into 308 closed cells and checking interval lower bounds. Thus $g$ is nonnegative on $[-T,T]$, with exactly the four zeros $\pm t,\pm T$. Every feasible measure of width at most $2T$, translated to that interval, must be supported on those zeros. Total mass and the first cosine equation fix the two pair totals, since $\cos(\lambda_1t)>0>\cos(\lambda_1T)$. The first two sine equations fix the differences within pairs to zero because the checked determinant
\[
 \sin(\lambda_1t)\sin(\lambda_2T)-\sin(\lambda_1T)\sin(\lambda_2t)
\]
is strictly negative. The candidate is therefore uniquely forced, and since it has both endpoint masses positive, $2T$ is the minimum. The interval on $2T$ gives the stated width bounds.

For reproducibility, all arithmetic assertions in the program use integers or exact rationals. An interval is represented by integer endpoints on a grid of step $10^{-30}$; multiplication and division by positive integers round outward. Trigonometric evaluations use the Taylor polynomial at a grid midpoint through degree 80 (cosine) or 81 (sine), with absolute remainder bounded by $10^{-30}$ on $|x|\le7$ via the factorial inequalities checked in the program. The interval is extended using the $1$-Lipschitz property of sine and cosine. Thus the finite checks are not a sampling argument or a floating-point optimization. Printed decimal summaries are rounded for readability and are not used in assertions. Running the script without disabling Python assertions reproduces the proof checks.
\end{proof}

\section{Consequences for the prime explicit formula}
For a smooth compactly supported weight $W$ on $(0,\infty)$ and a primitive Dirichlet character $\chi$, write $S_\chi(x)=\sum_{n\ge1}\Lambda(n)\chi(n)W(n/x)$, and $\widetilde W(s)=\int_0^\infty W(y)y^{s-1}\dd y$. In the classical smoothed explicit formula (by Mellin inversion and shifting the $-L'/L$ contour), each nontrivial zero $\rho$ of multiplicity $m_\rho$ contributes $-m_\rho\widetilde W(\rho)x^\rho$. The pole and trivial-zero terms are separate. Under the supplied DGRH, all nontrivial zeros are $\rho=1/2+i\gamma$. Because $W$ is smooth and compactly supported, its Mellin transform decays faster than any power in fixed vertical strips; the zero count $O(T\log(q(T+2)))$ permits termwise positive averaging over a fixed compact dilation interval.

Replacing $S_\chi(x)$ by $\int S_\chi(x\e^u)\dd\mu(u)$ multiplies each zero term by $F_\mu(\rho)$. Requiring a root of order $m$ also annihilates $x^\rho$ times any polynomial in $\log x$ of degree below $m$, since
\[
 \int (\log x+u)^k\e^{\rho u}\dd\mu(u)=0\quad(k<m).
\]
This is useful for uncertainty in zero locations and for differentiated explicit formulas. It is not necessary to impose $m=m_\rho$ merely to cancel an ordinary multiplicity-weighted zero term: simple vanishing already suffices. The additional jets enforce robustness. Taylor's theorem gives, for a filter on $[0,L]$ with order $m$ at $\rho$,
\[
 |F_\mu(\rho+\eta)|\le \e^{(\Rea\rho+|\eta|)_+L}\frac{(|\eta|L)^m}{m!}.
\]
Positive averaging does not cancel a real pole or a central real zero with positive exponential multiplier; the method targets nonzero ordinates.

For a common real part $\beta$, exponential tilting is a bijection between feasible measures for the complex-exponent and Fourier problems, preserving support diameter. Given an optimal Fourier measure $\nu$, the raw dilation filter is
\[
 \dd\mu(u)=\frac{\e^{-\beta u}\dd\nu(u)}{\int\e^{-\beta v}\dd\nu(v)}.
\]
The same orders of vanishing hold. In normalized prime errors the Fourier filter can be applied directly after extracting the factor $x^{1/2}$. Thus the minimum width proved above means an exact minimum possible ratio of largest to smallest dilation scale, $\e^D$, for positive filters satisfying the prescribed spectral constraints, irrespective of the prime weight. This is optimal for universal cancellation of those zero terms, not necessarily for a particular $W$ whose Mellin transform may already vanish.

\subsection{First two and first three zeta ordinates}
The reported first two ordinates have ratio $<3$. Corollary~\ref{cor:pair} gives the exact width $4\pi m/(\gamma_1+\gamma_2)$ for every equal jet order $m$. For the first three ordinates, use Theorem~\ref{thm:actual} and tilt at $\beta=1/2$. The minimum logarithmic window is about $0.469945213$ and the multiplicative ratio is about $1.59990654$. For comparison, using three separate simple two-atom filters would require width $\pi\sum_{i=1}^3 1/\gamma_i\approx0.497$, a strictly larger interval. (The individual pair optimum plus a separate third filter is another feasible benchmark.) The centered Fourier masses are approximately $0.15280708$ at each endpoint and $0.34719292$ at each inner node; the raw dilation weights are obtained by the exact tilt above, then by translation if desired.
\begin{center}
\begin{tabular}{lrr}
\toprule
Constraint & log width (approx.) & scale ratio (approx.)\\
\midrule
First two, order 1 at each & 0.357438197 & 1.429662206\\
First two, order 2 at each & 0.714876394 & 2.043934024\\
First two, order 3 at each & 1.072314591 & 2.922135226\\
First three, simple & 0.469945213 & 1.599906536\\
\bottomrule
\end{tabular}
\end{center}
The analytic pair theorem and the uniform rational certificate justify the optimality statements. The displayed decimals use high-precision numerical evaluations and are not all interval-certified to their printed digits. In particular the rigorously certified enclosure for the three-frequency width is the one in Theorem~\ref{thm:actual}, given the frequency boxes. The method applies to any primitive Dirichlet family using DGRH, without assumptions on zero spacing or independence; repeated targets must be combined before applying Theorem~\ref{thm:base} or Theorem~\ref{thm:perturb}.

\section{Limitations and further research}
\begin{enumerate}
\item The local analytic theorem does not classify global transitions for arbitrary three or more frequencies. At a base polynomial with zero coefficients, the Jacobian used above can lose rank, and the optimal number of atoms can change; the fractional powers seen in the article's two-frequency cusp suggest a singular perturbation problem.
\item Equal jet orders at two frequencies in the band $1<\gamma/\alpha\le3$ are settled, but arbitrary unequal orders and ratios outside the band remain open here. More general coefficient positivity regions for reflected spectra, including very wide bands rescued by extra central roots, are natural algebraic problems.
\item A quantitative perturbation radius could be extracted from a lower bound for the base Jacobian's smallest singular value, node separation and derivative margins. The certificate in Section~\ref{sec:cert} is one finite implementation of the positivity strategy, not a general uniform complexity estimate.
\item For many zeros, the article's $O(\log B)$ support problem remains: improve its leading constant and the Jensen lower bound. It is unclear whether near-reflection pairings can systematically improve global constructions.
\item A formalization in Lean could separate (i) the transform/translation and polynomial lemmas, (ii) the residue identity and positivity, (iii) the finite-dimensional analytic implicit function and normal-vector argument, and (iv) the rational interval certificate. The DGRH application would import the vendored \texttt{Main} in its own pinned Lake environment. This note is a mathematical proof and executed Python certificate, not a Lean-checked extension. No changes to the repository were made.
\end{enumerate}

\begin{thebibliography}{9}
\bibitem{vendor} ProveIt, vendored GRH development, read at main snapshot above. \href{https://github.com/VladimirReshetnikov/ProveIt/tree/a4404ca0247c31b8760589015b7eb6ab4efb6f92/Analysis/GeneralizedRiemannHypothesis/vendor/grh-lean}{Vendored source}. In particular \texttt{grh/src/Statements.lean}, \texttt{grh/src/Main.lean}, and \texttt{grh/README.md}.
\bibitem{article} ProveIt, \emph{Beyond the Riemann hypothesis for Dirichlet L-functions: consequences, transfers and limits}. \href{https://github.com/VladimirReshetnikov/ProveIt/tree/a4404ca0247c31b8760589015b7eb6ab4efb6f92/Analysis/GeneralizedRiemannHypothesis/Article}{Merged article}, especially \texttt{sections/src.tex}, \texttt{sections/rec.tex}, theorem labels \texttt{rec:thm:hblock}, \texttt{rec:thm:twofreq}, \texttt{rec:thm:cusp}, and the final open questions.
\bibitem{lau} M. A. Lau and L. Y. Pao, Input shaping and time-optimal control of flexible structures, \emph{Automatica} 39 (2003), 893--900. \href{https://doi.org/10.1016/S0005-1098(03)00024-4}{DOI}.
\bibitem{pao} Discrete time-optimal command shaping, \emph{Automatica} (2007), \href{https://www.sciencedirect.com/science/article/abs/pii/S0005109807001264}{publisher abstract and overview}; related work on derivative constraints and time-optimal discrete designs.
\bibitem{flint} FLINT documentation, \emph{The Riemann zeta function, its zeros and L-functions}, \href{https://flintlib.org/doc/tutorials/numerics/zeta.html}{numerical tutorial}, example \texttt{zeta\_zeros -n 1 -count 5 -digits 30}. Source of the reported ordinate digits; the zero-location computations themselves were not replayed.
\end{thebibliography}
\end{document}
